\begin{remarks}[6.15.11]
\label{IV.2.s6.15.11-fr}
(i) En la demostración de (6.15.10) no puede prescindirse
de (6.14.1), ni siquiera cuando $X=Y$, porque interviene
la clausura integral del anillo $B$, que no es necesariamente
noetheriano aunque $B$ lo sea.

(ii) El ejemplo (6.5.5, (ii)) muestra que, en el enunciado
de (6.15.10), no se puede reemplazar «geométricamente
unirramoso» por «íntegro», aunque el cuerpo residual $\boldsymbol{k}(x)$
sea algebraicamente cerrado y el morfismo $g$ sea étale.

\oldpage[IV-2]{182}
Tampoco puede reemplazarse en ese enunciado
«geométricamente unirramoso» por «unirramoso». Sea, en
efecto, $A$ el anillo local íntegro completo
$\mathbb R[[U,V]]/(U^2+V^2)$;
si $u$ y $v$ son las imágenes de $U$ y $V$ en $A$, el ideal
maximal de $A$ es $Au+Av$. Se comprueba fácilmente que la
clausura integral de $A$ es el anillo $A[t]$, donde $t=u/v$
satisface la relación $t^2=\setmin 1$, de modo que $A[t]$ es
isomorfo al anillo local $\mathbb C[[U]]$; como el cuerpo
residual de $\mathbb C[[U]]$ es $\mathbb C$, se ve que $A$
es unirramoso pero no geométricamente unirramoso. Sin
embargo, $A\otimes_{\mathbb R}\mathbb C$ no es un anillo
íntegro, pues es isomorfo a $\mathbb C[[U,V]]/(U\setmin iV)(U+iV)$.
\end{remarks}
